\section{Introduction}
In 2024, total global assets under management exceeded USD 128 Trillion, and some estimates projected that 60\% of this amount, or USD 87.6 Trillion, would be invested in active management by the end of 2025.\footnote{\url{https://web-assets.bcg.com/cc/0a/25876ea740168e908a8652e147d7/2025-gam-report-april-2025.pdf}
\url{https://www.pwc.com/ng/en/press-room/global-assets-under-management-set-to-rise.html}} These investments  differ from direct investments in companies, since investors have a much lower degree of influence on managers' decisions \citep{BebchukCohenHirst2017}. Asset managers usually determine their investment strategies through an advisory board in which investors have no or, at best, little voting power. Moreover, higher fees in the asset management industry are not necessarily a stamp of quality, as worse performing funds may charge higher fees
\citep{Gil-Bazo2008}.

This environment leads to clear incentives for managers to take decisions aimed towards increasing their own profits, without giving regard to the impact that they could have on investors' welfare. Moreover, since investors as a group tend to chase recent performance, and the relation between recent performance and assets under management is usually increasing and convex \citep{Ippolito1992ConsumerReaction, chevalier1997risk, sirri1998costly}, fund managers have an option-like payoff in their fee revenues. Therefore, assuming that fund managers have positive expected information ratios, one way to improve their profits is by increasing the volatility of the underlying portfolio \citep{BrownHarlowStarks1996}.
How should investors compensate managers to avoid taking excessive risk? In this paper we develop a principal-agent model of the asset management industry to answer this question. Moreover, we provide a simple method to provide the fair fee level as a linear function of the investor's desired volatility. We provide a numerical illustration with realistic parameters to show that this linear contract is a good approximation of the optimal contract. Our main result suggests that a properly chosen fee schedule can align incentives and implement a desired risk target.

We contribute to the literature on agency problems in delegated portfolio management. Broadly, delegated portfolio management is a classic principal--agent setting in which investors typically contract on outcomes (returns relative to benchmarks) rather than on the manager's underlying actions, because effort and information acquisition are not directly observable \citep{Holmstrom1979,AdmatiPfleiderer1997}. This creates scope for agency costs: managers may under-invest in costly research/monitoring when the effort is hard to verify \citep{Holmstrom1979,JensenMeckling1976}, engage in short-horizon appearance management around disclosure and evaluation dates \citep{Musto1999,CarhartKanielMustoReed2002}, or choose actions that raise their private benefits (fees, bonuses, career concerns) even when they do not maximize investors' welfare \citep{JensenMeckling1976}. These frictions motivate the design of incentive contracts in asset management--including performance-based fees and benchmark-linked pay---as studied in agency models of portfolio management and dynamic incentives \citep{Starks1987,AdmatiPfleiderer1997,HolmstromMilgrom1987}.
